% Included after the bibliography by the main report; do not compile independently.
\appendix[Individual Contributions and Reflections]
\label{app:personal-contributions}

This appendix records individual contributions to FallGuard. Yaorong Huang's entry describes his code, writing, and web-presentation responsibilities. The other entries remain TBD and must be completed by the corresponding members before submission. Author order does not imply a verified division of work.

\subsection{Yaorong Huang}
\subsubsection*{Code implementation and integration}
I carried out most of the project's code implementation and organization, covering temporal modeling, pose preprocessing, training/evaluation tools, model export, and edge--cloud integration. MaskedBiMamba combines established Mamba operators and pose estimation with temporal attention, bidirectional processing, frame masking, and quality weighting. My contribution concerns their composition, implementation, evaluation tools, and integration; the underlying operators and pose extractor originate from the cited work.

Public evidence includes \path{code/src/fallbench/} and \path{code/scripts/export_mamba_onnx.py}. The released \path{deployment/} directory contains the offline runtime, independent outbox, cloud puller, dashboards, and generic service templates. Selected temporal weights and checksums are provided in \path{artifacts/models/}.

\subsubsection*{Offline deployment and workspaces}
I integrated USB-camera/video acquisition, Hailo pose extraction with a CPU fallback, and the two-input MaskedBiMamba ONNX classifier on the Pi. Separating inference from the SQLite outbox preserves committed records during disconnection or camera stoppage. Cloud synchronization commits records before acknowledging UUIDs and retains original capture times during retry.

I completed the three-page local and cloud workspaces: live preview and fresh results, historical charts with filters/pagination/CSV, and device diagnostics. The earlier minimal local page provided insufficient independent observation and troubleshooting. The updated interface clears current probability for missing poses, expired results, or lost connections. The cloud retains its archive while current Pi status uses a protected private bridge; controls validate same-origin JSON and a page token.

Deployment verification preserved existing databases, checked service health and source hashes, and saved rollback backups. Temporary integration tests cover archive queries, restart persistence, authentication, delivery, and actual service startup/shutdown; both actual page scripts were checked for stale and invalid states. Public deployment evidence is in \path{results/deployment/}.

\subsubsection*{Technical issues and evidence boundaries}
Archived missing-joint tests zero-filled coordinates before normalization, also rescaling observed joints. I corrected the order while excluding the affected results from current conclusions; the fix does not validate those old experiments. I also separated window/video predictions, merged events, and fixed-rule confirmed alerts. Pipeline throughput, detection timing, cloud receipt latency, and operational archive counts describe different quantities and do not establish the same performance guarantee.

\subsubsection*{Manuscript and evidence organization}
I took primary responsibility for preparing and revising the LaTeX manuscript, method notation, protocols, results, implementation, and limitations. I maintained consistency across figures, aggregate evidence, technical explanations, deployed services, and repository materials. The interpretation retains TE's higher clean F1, mixed ablations, limited external discrimination, and incomplete fall-interval detection. Evidence is retained in \path{manuscripts/final-report/} and \path{results/}; external algorithms and datasets are credited.

\subsubsection*{Presentation and publication}
I developed the twelve-slide GitHub Pages presentation in \path{docs/}, including interactive charts, bilingual notes, searchable results, projection controls, a presenter console, and rehearsal timing. I synchronized the final report, Chinese guides, presentation, deployment source, and GitHub publication with the October workspace update. The presentation exposes saved evidence; the Pi and cloud provide the separate live demonstration.

\subsubsection*{Reflection}
A working prototype and a supported performance claim need different evidence. High clip recall can coexist with late detection; high video F1 can coexist with low specificity. Preprocessing can invalidate a robustness comparison. Deployment also requires freshness checks, independent storage, recovery, diagnostics, and consistent documentation. Future work should evaluate continuous recordings, calibrate alerts, and test longer camera operation across conditions, while preserving one traceable evidence chain.

\subsection{Jiahao Zhao}
TBD.
\subsection{Chengyi Xu}
TBD.
\subsection{Zhonglin Chang}
TBD.
\subsection{Jiayun Song}
TBD.
